\section{Discussion}
\label{sec:discussion}

\subsection{Clinical Implications}
\label{sec:clinical}

The predicted $n = 2$ frequency range (\SIrange{12}{18}{\hertz}) has
immediate relevance to occupational health.  Vehicle cabins expose operators
to broadband vibration spanning \SIrange{1}{20}{\hertz}, with dominant energy
typically concentrated below \SI{10}{\hertz} from engine and road
excitation~\cite{griffin1990handbook,mansfield2005human}.  The bladder's
fundamental mode sits at the upper edge of this band.  Two coupling mechanisms
are therefore relevant:

\paragraph{Sub-resonant forced response.}
At the ISO~2631 pelvic resonance peak (\SIrange{4}{8}{\hertz}), the bladder
responds as a forced sub-resonant oscillator.  The pelvic transmissibility is
at its maximum ($T \approx 2.3$ at \SI{5.5}{\hertz}), and while the bladder
modal amplification is below unity at these frequencies, the product
$T \times H$ remains significant.  This is the regime most relevant to
occupational WBV exposure: the bladder wall undergoes forced cyclic
deformation driven by amplified pelvic motion, even though the excitation
frequency is below the bladder's own resonance.

\paragraph{Near-resonant amplification.}
At the bladder's natural frequency (\SIrange{12}{18}{\hertz}), the modal
amplification reaches its peak ($Q = 2.5$), but pelvic transmissibility has
rolled off ($T \approx 0.3$).  This regime may be relevant to certain
industrial vibration sources (e.g., pneumatic tools, vibratory compactors)
that produce significant energy above \SI{10}{\hertz}.  The
fill-volume-dependent nature of the resonance means that the bladder's
susceptibility to a particular frequency shifts as it fills---a dynamic
vulnerability window.

\paragraph{Stretch receptor activation.}
Urgency sensation requires activation of stretch receptors
(mechanosensitive afferent nerves) in the detrusor muscle and
suburothelium~\cite{coste2010piezo1}.  The critical question is whether
vibration-induced cyclic wall displacement reaches the activation threshold
of these receptors.  Our model provides the wall displacement per unit
input; coupling this to afferent nerve thresholds---which are not
well-characterised in the dynamic regime---remains an open problem.


\subsection{Comparison with Other Pelvic Organs}
\label{sec:pelvic}

The bladder is not the only pelvic organ exposed to WBV.  The prostate,
uterus, and rectum are all fluid-containing or fluid-adjacent structures
embedded in the pelvic floor.  The present framework could, in principle,
be adapted to any of these organs by substituting appropriate geometric
and material parameters.  The bladder is simply the cleanest example: a
single-chamber, fluid-filled, roughly spherical shell with well-characterised
mechanical properties.


\subsection{Therapeutic Whole-Body Vibration}
\label{sec:therapeutic}

WBV at \SIrange{20}{40}{\hertz} is increasingly used for pelvic floor
rehabilitation in post-prostatectomy and postpartum
patients~\cite{Pelvic2019}.  The bio\-electrical activity of pelvic floor
muscles increases during synchronous WBV~\cite{Lauper2009}, and clinical
trials report improvements in continence
symptoms~\cite{NCT03325660}.  Our model suggests that these therapeutic
frequencies (\SIrange{20}{40}{\hertz}) lie above the bladder's $n = 2$ mode
for most fill states, placing the bladder in a supra-resonant regime where
wall displacement decreases with frequency.  The therapeutic effect is
therefore more likely mediated by pelvic floor muscle activation than by
direct bladder wall resonance.


\subsection{Limitations}
\label{sec:limitations}

Several simplifications warrant discussion.

\paragraph{Spherical geometry.}
The real bladder is slightly oblate and irregularly shaped, particularly at
low fill states.  A more realistic geometry would modify the mode shapes and
could split the degenerate frequencies of the spherical model.  However, at
moderate-to-high fill the spherical approximation is reasonable, and the
errors are of order $(1 - c/a)^2$ for the eigenfrequencies.

\paragraph{Pelvic constraint.}
Approximately 40\% of the bladder surface contacts the pelvic bone and
pelvic floor musculature.  A partial rigid boundary would modify the
effective mode shapes and could lower the resonant frequency by restricting
the free surface area.  This effect is not captured in the free-shell model
and represents the most significant limitation of the present analysis.

\paragraph{Nonlinear elasticity.}
The exponential strain-stiffening model (equation~\ref{eq:E_fill}) is a
phenomenological fit to vibrometry data.  A hyperelastic constitutive model
(Mooney--Rivlin or Ogden) would be more physically appropriate for large
deformations, though for the small-amplitude vibrations considered here,
a linearised modulus is adequate.

\paragraph{Layered wall structure.}
The bladder wall comprises distinct layers (mucosa, submucosa, detrusor
muscle, serosa) with different mechanical properties.  The present model
uses a single effective modulus.  A multilayer shell model would capture
the bending--stretching coupling between layers, potentially modifying
both frequencies and damping.

\paragraph{Surrounding tissue.}
The bladder is embedded in the pelvic cavity, surrounded by fat, connective
tissue, and adjacent organs.  This surrounding medium would add both
damping and mass loading, likely shifting the resonant frequencies downward
and broadening the peaks.  The free-shell model therefore represents an
upper bound on the resonant frequencies.
